\section{Section-by-Section Structural Notes}

\subsection*{Main paper}

\paragraph{Introduction and informal overview.}
The introduction does two things at once: it motivates the optimization question for two-layer networks, and it announces the PDE replacement of SGD as the central simplification. The ``informal overview'' is structurally important because it already gives the key formulas for $R_N$, $R(\rho)$, and the DD before the theorems appear.

\paragraph{Concrete examples.}
The isotropic Gaussian, anisotropic Gaussian, ReLU, and failure examples are not side material. They serve as demonstrations that the PDE viewpoint can both predict success and diagnose failure. The examples are also where the paper conveys intuition more clearly than in the generic theorems.

\paragraph{General results.}
This is the abstract core of the paper. Proposition~1 and Theorems~3--7 define the reusable mathematical package. If one only wants the proof architecture of the paper, this is the essential section in the main text.

\paragraph{Discussion and future directions.}
The discussion is modestly honest about what has and has not been proved. The PDE limit is strong, but the generic rates are incomplete, and the nicest convergence results require noisy / regularized dynamics. This section is useful because it tells the reader what not to overclaim.

\subsection*{Supplementary Information}

\paragraph{Section 2: static properties of $R(\rho)$.}
This section contains the mean-field variational backbone behind Proposition~1 and Corollary~2.1. It explains why minimizers can become singular and why the measure-valued perspective is nontrivial even though $R(\rho)$ looks convex-ish.

\paragraph{Section 3: dynamics and proof of Theorem 3.}
This is the most important supplement section mathematically. It contains the nonlinear dynamics representation, the continuity properties, and the trajectory-comparison lemmas that make the propagation-of-chaos argument work.

\paragraph{Sections 4 and 5: example-specific reduced analyses.}
These sections specialize the abstract machinery to the isotropic / anisotropic examples and related reduced PDEs. They are where one sees the practical analytical benefit of the mean-field reduction.

\paragraph{Section 6: noisy dynamics and free energy.}
This section contains the details behind the free-energy minimization, fixed points, and the noisy-DD convergence results. It is the technical foundation for Theorems~4--5.

\paragraph{Structural reading advice.}
If studying the paper for reuse rather than for every proof detail, the best order is:
\begin{enumerate}[label=\arabic*.]
\item read the informal overview for the main objects;
\item read the statements of Proposition~1 and Theorems~3--7;
\item read Supplementary Section~3;
\item read Supplementary Section~6;
\item only then revisit the example sections.
\end{enumerate}
